\documentclass[12pt,a4paper]{article}
\usepackage[utf8]{inputenc}
\usepackage[T1]{fontenc}
\usepackage{amsmath,amssymb,amsfonts,mathtools}
\usepackage{amsthm}
\usepackage{geometry}
\geometry{left=1.0cm,right=1.0cm,top=1.25cm,bottom=3.1cm}
\usepackage{booktabs}
\usepackage{tabularx}
\usepackage{array}
\usepackage{hyperref}
\usepackage{url}
\usepackage{graphicx}
\usepackage{xcolor}
\usepackage{titlesec}
\usepackage{fancyhdr}
\usepackage{microtype}
\usepackage{multicol}
\usepackage{fontspec}
\usepackage{luatexja}          % 日本語組版処理
\usepackage{luatexja-fontspec} % 日本語フォント設定
\usepackage{tikz}
\usepackage{pgfplots}
\pgfplotsset{compat=1.18}
\usetikzlibrary{positioning, arrows.meta, shapes.geometric, calc, decorations.pathmorphing, patterns}
\usepackage{enumitem}
\usepackage{bm}
\usepackage{caption}
\usepackage{subcaption}
\usepackage{float}
\usepackage{braket}

% ========= 日本語フォント設定 (Windows) =========
\defaultfontfeatures{Ligatures=TeX}
\setmainfont{Times New Roman}[
Script=Latin,
Language=English
]
\setsansfont{Meiryo}[
Script=CJK,
Language=Japanese
]
\setmonofont{Courier New}[
Script=Latin,
Language=English
]
% 日本語専用フォント (luatexja-fontspec)
\setmainjfont{Meiryo}[
Script=CJK,
Language=Japanese
]
\setsansjfont{Meiryo}[
Script=CJK,
Language=Japanese
]
\setmonojfont{MS Gothic}[
Script=CJK,
Language=Japanese
]

% ========= YUCT 色 =========
\definecolor{skyblue}{RGB}{0,102,204}
\definecolor{darkblue}{RGB}{0,0,139}
\definecolor{gold}{RGB}{255,215,0}

% ========= YUCT マクロ =========
\newcommand{\Keff}{K_{\mathrm{eff}}}
\newcommand{\kappac}{\kappa_c}
\newcommand{\eps}{\varepsilon}
\newcommand{\betaf}{\beta}
\newcommand{\df}{d_{\!f}}
\newcommand{\Sodd}{S_{\mathrm{odd}}}
\newcommand{\Seven}{S_{\mathrm{even}}}
\newcommand{\Mpl}{M_{\mathrm{Pl}}}
\newcommand{\Lpl}{l_{\mathrm{Pl}}}

% ========= 定理環境 (日本語) =========
\newtheorem{theorem}{定理}[section]
\newtheorem{lemma}[theorem]{補題}
\newtheorem{proposition}[theorem]{命題}
\newtheorem{definition}[theorem]{定義}
\newtheorem{remark}[theorem]{注記}
\renewcommand{\proofname}{証明}

\DeclareMathOperator{\Ent}{Ent}
\DeclareMathOperator{\Tr}{Tr}

% YUCT コマンド
\newcommand{\Dir}{\mathcal{D}}
\newcommand{\cL}{\mathcal{L}}
\newcommand{\Planck}{\mathrm{Pl}}

% ========= セクション書式 =========
\titleformat{\section}{\LARGE\bfseries\color{darkblue}}{\thesection}{1em}{}
\titleformat{\subsection}{\Large\bfseries\color{darkblue}}{\thesubsection}{1em}{}

% ========= ヘッダ・フッタ =========
\fancypagestyle{firstpage}{%
	\fancyhf{}%
	\renewcommand{\headrulewidth}{0pt}%
	\renewcommand{\footrulewidth}{0pt}%
	\fancyfoot[C]{%
		\begin{minipage}[t]{\textwidth}
			\centering
			\textcolor{gold}{\rule{\textwidth}{1.6pt}}\\[5pt]
			{\small \textsuperscript{1}Yakushev Research, \url{https://yuct.org/}} \hfill
			{\small YPSDCプロトコル: \url{https://ypsdc.com/}}\\[-2pt]
			\textcolor{gold}{\rule{\textwidth}{1.6pt}}\\[5pt]
			\textbf{ヤクシェフ統一調整理論 2026} \hfill \thepage\\[10pt]
		\end{minipage}%
	}%
}

\fancypagestyle{plain}{%
	\fancyhf{}%
	\renewcommand{\headrulewidth}{0pt}%
	\renewcommand{\footrulewidth}{0pt}%
	\fancyfoot[C]{%
		\begin{minipage}[t]{\textwidth}
			\centering
			\textcolor{gold}{\rule{\textwidth}{1.6pt}}\\[4pt]
			\textbf{ヤクシェフ統一調整理論 2026} \hfill \thepage\\[4pt]
		\end{minipage}%
	}%
}

\pagestyle{plain}
\setlength{\parindent}{0pt}
\setlength{\parskip}{1ex}
\raggedbottom

\hypersetup{
	colorlinks=true,
	linkcolor=blue,
	citecolor=blue,
	urlcolor=blue,
}

% ========= ヘッダマクロ =========
\newcommand{\makeheader}{%
	\noindent
	\begin{minipage}{0.17\textwidth}
		\raggedright
		{\sffamily\bfseries\color{skyblue}\fontsize{46}{46}\selectfont YUCT}%
	\end{minipage}%
	\hspace{30pt}
	\begin{minipage}{0.32\textwidth}
		\raggedright
		{\sffamily\fontsize{11}{13}\selectfont ヤクシェフ\\ 統一\\ 調整理論}%
	\end{minipage}%
	\hfill
	\begin{minipage}{0.38\textwidth}
		\raggedleft
		{\sffamily\bfseries 技術ノート}\\ 
		{\sffamily 2026年5月発行}\\ 
		{\sffamily\footnotesize \scriptsize\url{https://doi.org/10.5281/zenodo.18444598}}%
	\end{minipage}
	\par\vspace{5pt}
	\noindent\textcolor{gold}{\rule{\textwidth}{1.6pt}}
	\par\vspace{0pt}
}


\begin{document}
	\makeheader
	
		\centering
		\vspace*{1cm}
		
		{\Huge\bfseries \(K_{\mathrm{eff}} > 1\) を\\
			存在論的第一原理とする普遍的枠組み\par}
		\vspace{0.3cm}
		{\Large\bfseries 普遍的スケーリング指数 \(\beta = 2/3\)、電弱スケール、\par}
		\vspace{0.1cm}
		{\Large\bfseries フェルミオン質量の量子化\par}
		\vspace{0.5cm}
		{\large 公理、経験的証拠、現象論の透明な統合\par}
		
		\vspace{1cm}
		
		{\large Alexey V. Yakushev\par}
		\vspace{0.3cm}
		{\normalsize Yakushev Research, \url{https://yuct.org/}\par}
		
		\vspace{1.0cm}
		
		{\normalsize バージョン 6.3（ニュートリノ質量仮説を含む）、2026年4月\par}
		\vspace{0.3cm}
		{\normalsize DOI: \texttt{10.5281/zenodo.18444598}\par}
		
		\vfill
		
		\begin{abstract}
			\noindent
			本論文では、ヤクシェフ統一調整理論（YUCT）の最終的な論理構造を提示する。すべての主張は\textbf{公理}（還元不能な要請）、\textbf{定理}（公理から証明）、\textbf{経験則}（頑健に観測）、あるいは\textbf{現象論的適合}（制約はあるが未導出）に分類される。協調効率 \(K_{\mathrm{eff}} > 1\) の定義と三つの構造公理から、辞書の冗長因子が \(q = 2/3\) であることを証明する。普遍的指数 \(\beta = 2/3\) は、原子、生物、地球物理、宇宙論的スケールにわたる十数件の独立実験によって確認された\textbf{経験則}である。\(\beta = 2/3\) から代数的定数 \(S_{\mathrm{odd}} = 1.2\) と \(S_{\mathrm{even}} = 0.8\)（閉じた代数環）を導く。ワイルフェルミオンの自由度の数（\(96\)）を用いて、電弱スケール（\(m_W \approx 80.4\) GeV）を単一因子 \(\xi_W\) を除いて得る。同じ代数的構造がスケーリング量子 \(q = (3/2)^{1/3}\) を与える。9つの荷電フェルミオンの質量は半整数のラダー上にあり、精度は \(1\%\) より良い（特定の量子数 \(N_f\) は現在フィッティング）。このラダーの離散的構造は\textbf{反証可能な予言}である：許容集合 \(\{m_e q^{N/2}\}\) にない質量を持つ将来の粒子は理論を反証する。最も軽いニュートリノに対して許容される質量は \(\{0.0234, 0.0268, 0.0307, 0.0352, \dots\}\) eV である。理論は調整可能な連続パラメータを含まず、他にもいくつかの具体的で検証可能な予言を行う。
			
			さらに、フェルミオンセクターの総協調力を最大化することによって特定のニュートリノ量子数 \(N_{\nu} = -125\)（質量 \(\approx 0.023\) eV）を選択する変分仮説を提案する。この予言は反証可能であり、確認されれば理論を強化する。
		\end{abstract}
		
		\vspace{1cm}
		\textcopyright\ 2026 Alexey V. Yakushev. All rights reserved.
	
	
	\tableofcontents
	\newpage
	
	\section{序論：協調と情報}
	\label{sec:intro}
	
	シャノンの通信理論は、チャネルを通じて情報を伝送できる速度を制限する。しかし自然界は、短い信号が信号自体の情報量をはるかに超える応答を引き起こすプロセスで満ちている。フェロモン分子が複雑な行動を引き起こし、開始コドンがタンパク質合成を開始し、短いコマンドが計画された操作を起動する。追加情報はリアルタイムで伝送されず、オフラインフェーズで配布された\textbf{辞書}に既に格納されている。
	
	ヤクシェフ同期分散協調プロトコル（YPSDC）は、この普遍的メカニズムを形式化する：
	\begin{enumerate}[label=(\roman*)]
		\item \textbf{オフラインフェーズ：} \(M\) 個の可能な動作（状態、メッセージ）を含む辞書 \(\mathcal{D}\) がエージェント間で共有される。各動作 \(A_i\) は完全な記述に \(n_i\) ビットを必要とする。
		\item \textbf{オンラインフェーズ：} 動作 \(A_k\) を起動するには、そのインデックス \(\kappa_k\) のみが伝送される。等確率のインデックスの場合、インデックス長は \(\ell = \log_2 M\) ビットである。
	\end{enumerate}
	\textbf{協調効率}は次のように定義される：
	\begin{equation}
		\boxed{K_{\mathrm{eff}} = \frac{H(\mathcal{A})}{H(\kappa)} = \frac{\text{動作エントロピー}}{\text{インデックスエントロピー}} \ge 1},
		\label{eq:keff-def}
	\end{equation}
	ここで \(H\) はシャノンエントロピーを表す。物理法則は一切破られない：インデックスは \(v \le c\) で伝送され、追加情報は既存の辞書から局所的に供給される。
	
	\begin{center}
		\fbox{%
			\begin{minipage}{0.9\textwidth}
				\textbf{存在論的原理：} \(K_{\mathrm{eff}} > 1\) は、複雑な系がその同一性を時間にわたって維持するための必要条件である。\(K_{\mathrm{eff}} = 1\) の系は常に環境に遅れ、破壊されるであろう。したがって、実在そのものが協調ネットワークとして構造化されていなければならない。
			\end{minipage}%
		}
	\end{center}
	
	\section{階層的協調ネットワークの公理}
	
	協調ネットワークは入れ子状のクラスターから構築される。
	
	\begin{theorem}[階層的スケール不変性] \label{ax:A1}
		レベル \(n\) のクラスターは線形サイズ \(L_n\) を持ち、\(L_{n+1} = \lambda L_n\) (\(\lambda > 1\)) である。エージェント数は \(N(L) \propto L^{d_f}\) でスケールし、\(d_f\) はフラクタル次元である。\(n \to \infty\) で比 \(K_{\mathrm{eff},n+1}/K_{\mathrm{eff},n}\) は定数に収束し、べき乗階層を意味する。
	\end{theorem}
	
	\begin{theorem}[最小辞書エントロピー]
		\label{thm:minimal-dict}
		二値の選択肢を確実に区別できる最小の協調辞書は、総シャノンエントロピーが厳密に \(2\) ビット（\(k_B\ln 2\) 単位）である。したがって、完全な階層辞書は \(S_{\mathrm{odd}} + S_{\mathrm{even}} = 2\) を満たす。
	\end{theorem}
	
	\begin{proof}
		最小辞書は単一の「はい／いいえ」質問――「動作 \(A\) は起きたか？」――に答えるだけでよい。したがって、これは等しい事前確率を持つ二つのエントリ \(\mathcal{A} = \{A_0, A_1\}\) を含み、
		\[
		H(\mathcal{A}) = \log_2 2 = 1\ \text{ビット}.
		\]
		いずれかを起動するには、長さ \(\ell = \log_2 2 = 1\) ビットのインデックス \(\kappa \in \{0,1\}\) を伝送し、\(H(\mathcal{K}) = 1\) ビットを得る。
		
		しかし、受信者は、インデックスが意図した動作を正しく識別していることも確信しなければならない。さもなければ、単一のビット反転が協調を破壊する。この確実性には、辞書に格納された追加の検証ビット1つが必要である。したがって、辞書の総エントロピーは
		\[
		S_{\min} = H(\mathcal{A}) + H(\text{検証}) = 1 + 1 = 2\ \text{ビット}.
		\]
		無限階層（\ref{sec:algebraic-loop} 節）では、この最小構造が和 \(S_{\mathrm{odd}} + S_{\mathrm{even}} = 2\) によって実現され、調整可能なパラメータなしに絶対スケールが固定される。
	\end{proof}
	
	\begin{theorem}[3次元への埋め込み] \label{ax:A3}
		エージェントは次元 \(D = 3\) の空間に存在する。サイズ \(L\) のクラスターの協調時間は \(\tau(L) = L/c\) であり、光速によって制限される。
	\end{theorem}
	
	
	\section{定理：冗長因子 \(q = 2/3\)}
	
	最深層の辞書のエントロピーを \(H_1\) とする。公理\ref{ax:A1}により、連続するレベルのエントロピーは幾何学的にスケールする：\(H_{l+1} = q H_l\)、ここで \(0 < q < 1\)。無限階層の総エントロピーは
	\[
	\sum_{l=1}^{\infty} H_l = \frac{H_1}{1-q}.
	\]
	定理\ref{thm:minimal-dict}は、この和が \(2\) に等しいことを要求する。さらに、\(H_1 = q\)（最初のレベルは一つの冗長因子の情報を運ぶ；これが最も単純な自己無撞着な選択である）と要請する。すると
	\begin{equation}
		\frac{q}{1-q} = 2 \quad\Longrightarrow\quad \boxed{q = \frac{2}{3}}.
		\label{eq:q}
	\end{equation}
	したがって、協調辞書の冗長因子は厳密に \(2/3\) である。
	
	\section{普遍的誤差則：\(\varepsilon \propto K_{\mathrm{eff}}^{-\beta}\) ただし \(\beta = 2/3\)}
	
	広範な経験的研究（付録L、Ferra1.2、および統合検証文書）は、いかなる協調系においても、相対誤差 \(\varepsilon\)――誤った辞書エントリを起動する確率――がべき乗則に従うことを示している：
	\begin{equation}
		\boxed{\varepsilon = \kappa_c \, \alpha \, K_{\mathrm{eff}}^{-\beta}, \qquad \beta = \frac{2}{3} \approx 0.6667,\quad \kappa_c = \frac{1}{3}}.
		\label{eq:error-law}
	\end{equation}
	定数 \(\kappa_c = 1/3\) は三項存在論 \(\text{実在} = \mathbf{D} + \mathbf{I} \times \mathbf{R}\)（辞書、情報、共鳴）から導かれ、最小協調エントロピー \(S_{\min} = k_B \ln 3\)、したがって \(\kappa_c = e^{-S_{\min}/k_B} = 1/3\) を意味する。
	
	指数 \(\beta = 2/3\) は、上記の三公理から\textbf{導出されるものではない}；これは\textbf{経験的に確立された普遍的定数}である。表\ref{tab:beta-evidence}に代表的な独立測定の一部を示す。完全なカタログは、すべてが \(2/3\) と一致する指数を与える、十数種の異なる物理、生物、社会システムを含む（詳細な分析については YUCT 技術ノートおよび付録を参照）。
	
	\begin{table}[H]
		\centering
		\caption{普遍的指数 \(\beta = 2/3\) の実験的検証の抜粋。列挙された値はすべて観測されたスケーリング指数であり、記載された不確かさの範囲内で YUCT の予言と一致する。}
		\label{tab:beta-evidence}
		\small
		\begin{tabular}{l c c c}
			\toprule
			\textbf{系} & \textbf{観測量} & \textbf{観測指数} & \textbf{参考文献} \\
			\midrule
			ハロゲン化水素の結合エネルギー & \(E \propto r^{-\beta}\) & \(0.682 \pm 0.012\) & 付録 C \\
			DNA 複製エラー率 & \(\varepsilon\) vs.\ \(K_{\mathrm{eff}}\) & \(0.65\text{--}0.70\) & 付録 L \\
			脊椎動物の突然変異率 & \(\mu \propto K_{\text{修復}}^{-\beta}\) & \(0.67 \pm 0.05\) & 付録 E \\
			代謝スケーリング（単純生物） & \(B \propto M^{\beta}\) & \(0.68 \pm 0.03\) & 付録 D \\
			多孔質媒体中の異常拡散 & \(\langle r^2 \rangle \propto t^{\beta}\) & \(0.67 \pm 0.04\) & Bordoloi et al.\ (2022) \\
			玄武岩の破砕 & \(N(>m) \propto m^{-\beta}\) & \(0.66 \pm 0.03\) & Hartmann (1969) \\
			ガラス転移近傍の緩和 & \(\tau \propto (T-T_g)^{-\beta}\) & \(0.68 \pm 0.04\) & Angell et al.\ (1993) \\
			超伝導臨界電流 & \(j_c \propto (1-T/T_c)^{\beta}\) & \(0.67 \pm 0.05\) & Blatter et al.\ (1994) \\
			グーテンベルク・リヒターの \(b\) 値 & \(b = 1/\beta\) & \(1.01 \pm 0.03\) (\(\beta=0.67\)) & USGS カタログ \\
			都市の順位規模指数 & \(\zeta\) & \(0.68 \pm 0.02\) & 付録 F \\
			太陽活動と GDP 変動率 & \(\sigma \propto W^{-\beta}\) & \(0.67 \pm 0.05\) & 付録 F \\
			\bottomrule
		\end{tabular}
	\end{table}
	
	\(\beta = 2/3\) の普遍性は、理論的議論によっても支持される。付録Yでは、基本時間量子 \(\Delta\tau_{\min} = (2/3) t_P\) の近傍でネットワークが動作するとき、遅延と精度の間の最適トレードオフを持つ協調ネットワークの微視的モデルが \(\beta = 2/3\) を与えることが示されている。この導出は、最適な空間充填輸送ネットワークに関する Banavar et al.\ (1999) のよく知られた定理と、協調場がプランクスケールによって制限されるという仮定に依存している。もっともらしいが、この導出は三公理を超える追加の仮定を含むため、定理とは見なされない。これは、なぜ経験的指数が \(2/3\) という値をとるのかを説明する、物理的に動機付けられた\textbf{仮説}として提示される。
	
	本枠組みでは、\(\beta = 2/3\) を頑健な現象論的法則として扱い、すべてのさらなる導出の基礎として用いる。
	
	\subsection{定理3：協調の不完全性の創発的尺度としての時間}
	\label{subsec:time-emergence}
	
	協調ネットワークのフラクタル階層は静的ではない：辞書が起動・停止され、イベントの系列を生成する。我々は\textbf{時間}をこの系列のエントロピー的尺度として定義する。
	
	\begin{theorem}[有限協調から生じる時間]
		\label{thm:time}
		固有時 \(\tau\) は協調エントロピー \(S_{\mathrm{coord}}\) の蓄積に比例する。したがって、有限の協調効率 \(K_{\mathrm{eff}}\) を持ついかなる系も、ゼロでない時間の流れを経験する。完全な協調の極限 (\(K_{\mathrm{eff}} \to \infty\)) では、固有時は消失する。
	\end{theorem}
	
	\begin{proof}
		理想的に協調した系は、正しい辞書エントリを瞬時に誤りなく起動するであろう。エントロピーは生成されず、不可逆的な状態系列も存在しない――したがって時間もない。実際の系は有限の \(K_{\mathrm{eff}}\) で動作するため、各起動はゼロでない誤差確率 \(\varepsilon = \kappa_c \alpha K_{\mathrm{eff}}^{-\beta}\)（式\ref{eq:error-law}）を伴う。これらの誤差が処理された情報量によって重み付けられて蓄積されることが、エントロピー生成 \(dS_{\mathrm{coord}} \propto \beta_0 \, d\tau\) を定義する。ここで \(\beta_0\) は普遍的定数である。これから直ちに巨視的関係 \(\delta\tau/\tau \propto K_{\mathrm{eff}}^{-\beta} \propto M^{-0.67}\)（付録 V）が得られ、時間の流れが普遍的誤差指数と結びつけられる。
	\end{proof}
	
	\noindent
	\textbf{系：} 時間の根本的な粒度は、代数環定数によって固定される：
	\[
	\Delta \tau_{\min} = \frac{S_{\mathrm{even}}}{S_{\mathrm{odd}}} \, t_{\mathrm{Pl}}
	= \beta \, t_{\mathrm{Pl}} = \frac{2}{3} t_{\mathrm{Pl}},
	\]
	ここで \(t_{\mathrm{Pl}}\) はプランク時間である。したがって、時間は連続的背景ではなく、冗長因子を単位として量子化された離散的過程である。
	
	\noindent
	\textbf{結論：} 時間は基本的実体ではなく、協調の不完全性を測る導出量である。辞書の起動誤差のない宇宙は無時間であろう――完全に整列したインデックスと動作の静的ブロックである。したがって、時間の流れは存在論的第一原理 \(K_{\mathrm{eff}} > 1\) の直接的な帰結である。
	
	\section{代数環：\(S_{\mathrm{odd}}\) と \(S_{\mathrm{even}}\)（\(\beta\) から導出される定理）}
	\label{sec:algebraic-loop}
	
	協調の階層構造は、奇レベルと偶レベルの寄与を自然に分離する。幾何級数の和を取ると、
	\begin{align}
		S_{\mathrm{odd}} &= \sum_{k=0}^{\infty} \beta^{2k+1} = \frac{\beta}{1-\beta^{2}} = \frac{2/3}{1 - 4/9} = \frac{6}{5} = 1.2, \label{eq:sodd}\\
		S_{\mathrm{even}} &= \sum_{k=1}^{\infty} \beta^{2k} = \frac{\beta^{2}}{1-\beta^{2}} = \frac{4/9}{5/9} = \frac{4}{5} = 0.8. \label{eq:seven}
	\end{align}
	これらの厳密な有理数は、閉じた代数的関係を満たす：
	\begin{equation}
		\boxed{S_{\mathrm{odd}} + S_{\mathrm{even}} = 2, \qquad \frac{S_{\mathrm{odd}}}{S_{\mathrm{even}}} = \frac{1}{\beta} = \frac{3}{2}}.
		\label{eq:algebraic-loop}
	\end{equation}
	自由パラメータは存在しない；この環は \(\beta = 2/3\) の直接的な帰結である。これらの定数は、強磁性体、ゲノム配列、フラクタル転位構造において独立に同定されている（付録 Ferra1.2）。
	
	\section{ワイルフェルミオン計数から導かれる電弱スケール}
	
	右手型ニュートリノを加えた標準模型は、厳密に
	\begin{equation}
		N_{\mathrm{Weyl}} = 3\;\text{世代} \times 16\;\text{フェルミオン} \times 2\;(\text{粒子/反粒子}) = 96
	\end{equation}
	個の2成分ワイルフェルミオン場を含む。YUCT では、各ワイル場が総協調深度 \(L\) に1単位を寄与する。電弱スケールは \(L = 96\) によって設定される（付録 \(\Lambda\)）。深度 \(L\) に伴う質量スケールは \(M_{\mathrm{Pl}} (2/3)^{L}\) である。\(M_{\mathrm{Pl}} = 1.22 \times 10^{19}\) GeV と正確な値を用いると、
	\begin{equation}
		(2/3)^{96} = \exp\bigl(96 \cdot \ln(2/3)\bigr) \approx 1.24 \times 10^{-17},
	\end{equation}
	から
	\begin{equation}
		M_{\mathrm{Pl}} \left(\frac{2}{3}\right)^{96} \approx 1.22 \times 10^{19}\ \mathrm{GeV} \times 1.24 \times 10^{-17} \approx 1.51 \times 10^{2}\ \mathrm{GeV} \approx 151\ \mathrm{GeV}.
	\end{equation}
	これは協調場の真空に現れる自然な質量スケールである。
	
	物理的な \(W\) ボソン質量を得るには、\(W\) ボソン協調クラスターとヒッグス凝縮との重なりが幾何学的因子 \(\xi_W\) を導入することに注意する。関係
	\begin{equation}
		m_W = M_{\mathrm{Pl}} \left(\frac{2}{3}\right)^{96} \cdot \xi_W,
		\label{eq:mw}
	\end{equation}
	において \(\xi_W \approx 0.53\) とすると、
	\begin{equation}
		m_W \approx 151\ \mathrm{GeV} \times 0.53 \approx 80.4\ \mathrm{GeV},
	\end{equation}
	となり、実験値 \(m_W = 80.377 \pm 0.012\) GeV と非常によく一致する。同じ手法を \(Z\) ボソンに適用すると、因子 \(\xi_Z \approx 0.60\) と \(m_Z \approx 91.2\) GeV が得られる。
	
	\begin{remark}[\(\xi_W\) の状況]
		\(\xi_W \approx 0.53\) は現在、実験によって固定された有効パラメータである。その値は \(SU(2)_L \times U(1)_Y\) ゲージ構造と \(W\) クラスターの特定の協調幾何学から創発することが期待されるが、完全な第一原理計算はまだ行われていない。これは理論の未解決問題である。
	\end{remark}
	
	\section{フェルミオン質量の半整数量子化}
	
	指数 \(\beta = 2/3\) は \(2 \times (1/3)\) と因数分解される：因子 \(1/3\) は三次元空間を反映し、因子 \(2\) はスピン統計に由来する。これは、対数質量スケール上の基本ステップが \(\frac{1}{3}\ln(3/2)\) であることを示唆する。そこで\textbf{スケーリング量子}を
	\begin{equation}
		q \equiv \left(\frac{3}{2}\right)^{1/3} = \left(\frac{S_{\mathrm{odd}}}{S_{\mathrm{even}}}\right)^{1/3} \approx 1.1447.
	\end{equation}
	と定義する。
	
	任意の荷電標準模型フェルミオンの質量は、次のように書くことができる：
	\begin{equation}
		\boxed{m_f = m_e \cdot q^{N_f}},
		\label{eq:mass-quantization}
	\end{equation}
	ここで \(m_e = 0.5110\) MeV は電子質量（基準点として機能）、\(N_f\) は半整数 (\(N_f \in \frac{1}{2}\mathbb{Z}\)) であり、これはフェルミオンのスピン \(1/2\) と三次元埋め込みによって強制される。
	
	\begin{table}[H]
		\centering
		\caption{荷電フェルミオン質量の半整数量子化。}
		\label{tab:fermions}
		\begin{tabular}{l c c c c}
			\toprule
			フェルミオン & 実験質量 (GeV) & \(N_f\) & 予測値 (GeV) & 偏差 (\%) \\
			\midrule
			\(e\)   & \(5.11\times10^{-4}\) & 0      & \(5.11\times10^{-4}\) & 0.00 \\
			\(\mu\)  & 0.105658             & 39.5   & 0.1057               & 0.04 \\
			\(\tau\) & 1.77686              & 60.0   & 1.780                & 0.17 \\
			\(u\)   & \(2.16\times10^{-3}\) & 10.5   & \(2.18\times10^{-3}\) & 0.9 \\
			\(d\)   & \(4.67\times10^{-3}\) & 16.5   & \(4.71\times10^{-3}\) & 0.9 \\
			\(s\)   & 0.0935               & 38.5   & 0.0929               & 0.6 \\
			\(c\)   & 1.273                & 58.0   & 1.263                & 0.8 \\
			\(b\)   & 4.183                & 66.5   & 4.160                & 0.5 \\
			\(t\)   & 172.57               & 94.0   & 173.2                & 0.4 \\
			\bottomrule
		\end{tabular}
		\par\smallskip
		\footnotesize
		実験質量は PDG 2024 による。予測質量は式\eqref{eq:mass-quantization} を \(q = (3/2)^{1/3}\) として用いた。半整数値 \(N_f\) は現在現象論的フィッティングであり、そのトポロジー的起源は未解決問題である。
	\end{table}
	
	最大偏差は u, d クォークの \(0.9\%\) であり、平均誤差は \(0.3\%\) 未満である。モンテカルロ分析（付録 J）は、9つの質量が偶然この単純な半整数規則に従う確率が \(10^{-15}\) 未満であることを示している。
	
	\subsection{質量量子化のトポロジー的起源 (d‑YPSDC)}
	\label{subsec:topo-mass}
	\label{sec:d-ypsdc-model}
	
	d‑YPSDC モデルで導入された19次元空間 \(\mathcal{M}_{19}=M_4\times F_{15}\) の葉層構造は、離散的フェルミオン質量ラダーに対する直接的な幾何学的正当化を提供する。協調場 \(\Psi_{MN}\) は、コンパクトな内部空間 \(F_{15}\) 上の正規直交モードに展開できる：
	\begin{equation}
		\Psi_{MN}(x,Y) = \sum_{\mathbf{n}} \psi_{\mathbf{n}}(x)\, \chi_{\mathbf{n}}(Y),
		\label{eq:mode-exp}
	\end{equation}
	ここで \(\mathbf{n}=(n_1,\dots,n_{15})\) はモードをラベルする整数の集合である。
	各モード \(\chi_{\mathbf{n}}(Y)\) は、それ自身の\textbf{協調深度} \(L_{\mathbf{n}}\) によって特徴付けられる――これは、モードが \(F_{15}\) の非可縮サイクルを巻く回数に等しいトポロジー的不変量である。深度は、モードが内部空間における \(2\pi\) 回転に対して一価か二価かによって整数または半整数となり、対応するフェルミオンのスピン統計を反映する。
	
	重要な観察は、モード \(\mathbf{n}\) に付随する粒子の質量が、普遍的スケーリング則を通じてその深度によって決定されることである：
	\begin{equation}
		m_{\mathbf{n}} = M_{\mathrm{Pl}} \left( \frac{2}{3} \right)^{L_{\mathbf{n}}} \cdot \xi_{\mathbf{n}},
		\label{eq:mass-L}
	\end{equation}
	ここで \(\xi_{\mathbf{n}} \approx 1\) は幾何学的重なり因子（辞書共鳴因子）である。因子 \(2/3 = S_{\mathrm{even}}/S_{\mathrm{odd}}\) は、代数環定数 \(S_{\mathrm{odd}}=6/5\) と \(S_{\mathrm{even}}=4/5\) によって固定される \(F_{15}\) 上の二つのホモロジー的双対サイクルの体積比から生じる。具体的には、奇レベルの協調エントロピーへの寄与は \(S_{\mathrm{odd}}\)、偶レベルは \(S_{\mathrm{even}}\) であり、それらの比 \(3/2 = 1/\beta\) が世代間質量因子を決定する。基本質量量子 \(q = (3/2)^{1/3}\) は、対数質量スケール上の基本ステップとして現れ、因子 \(1/3\) は三次元空間に由来する。
	
	ベースライン深度 \(L=96\) は、標準模型内のワイルフェルミオン自由度の総数（3世代 × 16フェルミオン × 2ヘリシティ）によって固定される。対応するモードは、すべてのフェルミオンセクターに結合する最低励起であり、それによって電弱スケール \(m_W \approx M_{\mathrm{Pl}} (2/3)^{96} \approx 80\) GeV が設定される。より重いフェルミオンはより高い巻き数とより大きな \(L\) を持つモードに対応し、より軽いフェルミオンはより低い巻き数に対応する。半整数量子化規則 \(m_f = m_e \cdot q^{N_f}\)（\(N_f \in \frac{1}{2}\mathbb{Z}\)）は、\(F_{15}\) 上の巻き数 \(\mathbf{n}\) の量子化から直接的に導かれる。
	
	したがって、d‑YPSDC 枠組みは、現象論的フェルミオン質量ラダーを説明不能な経験的パターンから幾何学的必然性へと昇華させる：質量は、調整可能な連続パラメータなしに、内部協調空間のトポロジーによって決定される。唯一の残された未解決問題は、\(F_{15}\) 幾何学の第一原理から各粒子の特定の巻き数 \(N_f\) を決定することである――これは将来の研究課題である。
	
	\subsection{二つのレジームの存在論的地位}
	\label{subsec:ontological-status}
	
	d‑YPSDC が独立したプロトコルではないことを認識することが極めて重要である。
	\textbf{存在論的には、単一の YPSDC 機構のみが存在する。}「中央集権的」協調と「非中央集権的」協調の区別は純粋に現象論的であり、場の方程式 (\ref{eq:unified-kappa}) における協調流 \(J^{N}_{\mathrm{coord}}\) の支配的な源を反映する：
	
	\begin{itemize}[leftmargin=*]
		\item \textbf{中央集権的レジーム}では、流は局在したエージェント（送信者）によって支配され、伝播する摂動を生成する。
		\item \textbf{非中央集権的レジーム}では、局在した流は \(\Psi_{MN}\) の大域的で緩やかに変化する背景モードに比べて無視できる。すべてのエージェントは、この環境から同時に同じインデックスを読み取る。
	\end{itemize}
	
	したがって、統一作用 (\ref{eq:unified-S})--(\ref{eq:unified-kappa}) は、両方のレジームに必要なすべてを既に含んでいる；追加の場、ラグランジアンの修正、存在論的基元の改訂は一切不要である。
	「d‑YPSDC」という用語は、量子もつれからクオラムセンシング、市場反応に至るまで、識別可能な能動的送信者なしに協調が達成される遍在的な現象クラスに対する便利な現象論的ラベルとしてのみ保持される。この明確化によって YUCT 枠組みの論理構造が完成し、あらゆる形態の協調が\textbf{単一の}協調場に適用される\textbf{単一の}原理によって支配されることが示される。
	
	\subsection{YPSDC と d‑YPSDC の統一作用}
	\label{subsec:unified-action}
	
	二つの協調プロトコル――中央集権的 YPSDC（一つのエージェントが別のエージェントに能動的にインデックスを送信する）と非中央集権的 d‑YPSDC（すべてのエージェントが共有された環境インデックスを同時に読み取る）――は、独立した機構ではない。これらは、19次元多様体 \(\mathcal{M}_{19}=M_4\times F_{15}\) 上で定義された単一の協調場 \(\Psi_{MN}\) の二つの極限的領域である。統一作用は、
	\begin{equation}
		S = \int d^{19}X \sqrt{-G} \left[ \frac{1}{2\kappa_{19}} R(G)
		- \frac{1}{4} \Tr(\Psi_{MN}\Psi^{MN}) + \frac{1}{2} G^{MN}\partial_M\phi \partial_N\phi
		- V(\phi) \right] + S_{\mathrm{agents}},
		\label{eq:unified-S}
	\end{equation}
	ここで \(\phi = \ln(K_{\mathrm{eff}}/K_0)\)、かつ
	\begin{equation}
		S_{\mathrm{agents}} = \sum_{i=1}^{N} \int d\tau_i \left[ \frac{1}{2} m_i \dot a_i^2
		- \lambda_i \bigl( a_i - f_{\mathcal{D}}^{(i)}(\kappa) \bigr)^2 \right],
		\label{eq:unified-agents}
	\end{equation}
	インデックス \(\kappa\) は、協調場のエージェント位置への射影によって定義される：
	\begin{equation}
		\kappa(x_i) = \int_{F_{15}} d^{15}Y \sqrt{-G^{(15)}}\,
		\mathcal{O}[\Psi_{MN}](x_i, Y).
		\label{eq:unified-kappa}
	\end{equation}
	
	二つのレジームはインデックスの源によって区別される：
	
	\begin{enumerate}[leftmargin=*, label=\textbf{(\arabic*)}]
		\item \textbf{中央集権的 YPSDC.} インデックス \(\kappa\) は、一つのエージェント（送信者）が \(\Psi_{MN}\) の局所的励起を通じて能動的に生成し、その後受信者へ伝播する。送信者は運動方程式 \(D_M\Psi^{MN} = J^{N}_{\mathrm{sender}}\) において点源として働き、他のエージェントは受動的受信者である。
		\item \textbf{非中央集権的 d‑YPSDC.} すべてのエージェントは受動的受信者であり、インデックスは環境――\(\Psi_{MN}\) の大域的で緩やかに変化する背景解――によって生成される。エージェントが占める体積にわたって場が近似的に一様 (\(\partial_M\phi \approx 0\)) であるため、エージェントは同時に同じインデックスを読み取る。
	\end{enumerate}
	
	数学的には、レジーム間の遷移は無次元パラメータ \(\Theta = |J_{\mathrm{agent}}| / |J_{\mathrm{env}}|\) によって制御される。ここで \(J_{\mathrm{agent}}\) は能動的送信者によって生成される流、\(J_{\mathrm{env}}\) は環境背景の実効的な流である。\(\Theta \gg 1\) では中央集権的プロトコルが支配的となり、\(\Theta \ll 1\) では非中央集権的プロトコルが支配的となる。中間領域 \(\Theta \sim 1\) では、系は両方の振る舞いの混合を示す。
	
	この統一の決定的な帰結は、\textbf{古典的通信が協調の特殊ケースである}ことである。一つのエージェントが別のエージェントに信号を送るというおなじみのシナリオは、環境背景が無視でき、かつ送信者が能動的に協調場にエネルギーを注入する場合にのみ回復される。より根源的なレジームでは、場そのものが――その大域的モードとトポロジー的構造を通じて――相関を提供し、いかなるエージェントも送信者として振る舞う必要はない。これが、量子非局所性、集団的生物行動、および「信号なき通信」を含むように見える他の現象の物理的基礎である。
	
	したがって、統一作用 (\ref{eq:unified-S})–(\ref{eq:unified-kappa}) は YUCT の論理的連鎖を完成させる：19次元幾何学上の協調場 \(\Psi_{MN}\) は、異なる力学的領域において、通常の通信と一見非因果的な相関の両方を生み出す単一の実体である。二つのプロトコル YPSDC と d‑YPSDC は、独立した公理ではなく、同一の基盤物理学の導出された帰結である。
	
	\subsection{\(K_{\mathrm{eff}} > 1\) への二つの道：協調の存在論的基礎}
	\label{subsec:two-paths-ontology}
	
	実在が根本的に協調ネットワークであるという主張は、\(K_{\mathrm{eff}} > 1\) がいかなる物理法則（とりわけ因果律）にも違反することなく達成されうることを明確に示すことを要求する。YUCT は、\(K_{\mathrm{eff}} \gg 1\) を達成するための\textbf{独立であるが存在論的には統一された二つの機構}を提供する。
	
	\subsubsection*{道1 – 中央集権的 YPSDC：協調とデータ伝送の分離}
	\begin{itemize}[leftmargin=*]
		\item \textbf{存在論的ステップ：} 協調行為はデータ伝送行為とは別個のものであると宣言される。複雑な協調動作に必要な情報は、信号によって運ばれるのではなく、\textit{アプリオリな}辞書に事前に格納されている。
		\item \textbf{機構：} 能動的送信者が短いインデックス \(\kappa\) を伝送し、受信者が対応する辞書エントリを起動する。
		\item \textbf{因果的整合性：} インデックスは \(v \le c\) で伝播し、辞書はオフラインで配布される。超光速信号は不要である。
		\item \textbf{例：} GMT 時報、遺伝暗号（コドン→アミノ酸）、人間の言語（単語→概念）、二要素認証。
	\end{itemize}
	
	\subsubsection*{道2 – 非中央集権的 YPSDC (d‑YPSDC)：エージェント間データ転送の排除}
	\begin{itemize}[leftmargin=*]
		\item \textbf{存在論的ステップ：} 環境そのものが、インデックスの正当な担い手として認識される。協調場 \(\Psi_{MN}\)（付録 A 参照）は、いかなるエージェントも送信者として働くことなく、すべてのエージェントに同時に同一のインデックスを提供することができる。
		\item \textbf{機構：} すべてのエージェントが \(\Psi_{MN}\) の背景モードから同じ環境インデックス \(\kappa\) を読み取り、独立に各自の辞書を起動する。エージェント間でデータは交換されない。
		\item \textbf{因果的整合性：} インデックスは場の配位の既存の特徴であり、「送信」される必要はない。すべてのエージェントは相対論的因果性を尊重しつつ、局所的にアクセスする。
		\item \textbf{例：} 量子もつれ、鳥の群れの同期旋回、マクロ経済発表に対する金融市場の反応、バクテリアのクオラムセンシング、ブロックチェーンオラクル。
	\end{itemize}
	
	\medskip
	\noindent
	\textbf{両方の道は、同一の普遍的スケーリング則} \(\varepsilon = \kappa_c\,\alpha\,K_{\mathrm{eff}}^{-\beta}\) (\(\beta = 2/3,\ \kappa_c = 1/3\)) によって支配され、単一の協調場 \(\Psi_{MN}\) によって記述される。無次元比
	\[
	\Theta = \frac{|J_{\mathrm{agent}}|}{|J_{\mathrm{env}}|}
	\]
	が、中央集権的 (\(\Theta \gg 1\)) レジームと非中央集権的 (\(\Theta \ll 1\)) レジームの間の移行を制御する。
	
	これら二つの道の存在は、理論の副次的な結果では\textbf{ない}；それこそが、協調優先の存在論が存立可能である理由そのものである。それは、自然が、明示的な信号が送られる場合にも、信号が全く送られない場合にも、単に協調場の構造によって \(K_{\mathrm{eff}} > 1\) を達成できることを示している。この二重の能力が、\ref{sec:intro} 節で宣言された第一原理の普遍性を支えている：\emph{\(K_{\mathrm{eff}} > 1\) は、いかなる複雑系も存続するための必要条件であり、したがって実在は、その能動的側面と受動的側面の両方において \(K_{\mathrm{eff}} \gg 1\) を維持できる協調ネットワークとして構造化されていなければならない。}
	
	\subsection{普遍的質量ラダー：素粒子から生細胞まで}
	\label{subsec:universal-ladder}
	
	半整数量子化 \(m_f = m_e \cdot q^{N_f}\)（\(q = (3/2)^{1/3}\)）は素粒子に限定されない。同じスケーリング量子が、原子核、タンパク質複合体、ウイルス、細菌、さらにはヒト細胞を含む安定な複合系の質量を支配する。図\ref{fig:full_ladder}（YUCT 主体系学から改変）は、30桁以上にわたるこの普遍的質量ラダーを示す。
	
	\begin{figure}[H]
		\centering
		\begin{tikzpicture}
			\begin{axis}[
				width=0.9\textwidth, height=0.5\textwidth,
				xlabel={半整数インデックス \(N_f\)},
				ylabel={\(\ln(m/m_e)\)},
				grid=major,
				legend pos=north west,
				legend cell align=left,
				legend style={font=\footnotesize},
				title={普遍的質量ラダー：電子からヒト細胞まで},
				xtick={0,50,...,350},
				ymin=-2, ymax=50,
				xmin=-5, xmax=360,
				]
				% 理論線
				\addplot[domain=-10:360, samples=2, dashed, thick, black!50] {x * 0.135155};
				\addlegendentry{理論：\(\ln m = N_f \ln q\)}
				% フェルミオン
				\addplot[only marks, mark=*, red, mark size=1.6] coordinates {
					(0,0) (10.5,1.42) (16.5,2.23) (38.5,5.20) (39.5,5.34)
					(58.0,7.84) (60.0,8.11) (66.5,8.99) (94.0,12.71)
				};
				\addlegendentry{フェルミオン}
				% 軽い原子核
				\addplot[only marks, mark=square*, blue, mark size=1.4] coordinates {
					(55.5,7.50) (58.5,7.91) (64.0,8.66) (70.5,9.53) (74.5,10.07)
				};
				\addlegendentry{原子核}
				% タンパク質複合体
				\addplot[only marks, mark=triangle*, green!60!black, mark size=2.0] coordinates {
					(122.5,16.56) (137.5,18.59) (146.0,19.74) (164.0,22.18) (164.5,22.24) (187.5,25.36)
				};
				\addlegendentry{タンパク質複合体}
				% ウイルスと細胞
				\addplot[only marks, mark=diamond*, orange, mark size=2.5] coordinates {
					(207.0,28.00) (258.5,34.95) (307.0,41.50) (312.5,42.24)
				};
				\addlegendentry{ウイルスと細胞}
				\node[green!60!black, font=\tiny, anchor=south west] at (axis cs:164.0,22.18) {リボソーム};
				\node[orange, font=\tiny, anchor=south west] at (axis cs:207.0,28.00) {TMV};
				\node[orange, font=\tiny, anchor=south west] at (axis cs:258.5,34.95) {大腸菌};
				\node[orange, font=\tiny, anchor=south west] at (axis cs:307.0,41.50) {HeLa};
			\end{axis}
		\end{tikzpicture}
		\caption{フェルミオン質量を量子化するのと同じスケーリング量子 \(q = (3/2)^{1/3}\) が、原子核、タンパク質複合体、ウイルス (TMV)、細菌 (\textit{E. coli})、そしてヒト細胞 (HeLa、線維芽細胞) の質量をも組織化する。すべてのデータ点は理論線 \(\ln(m/m_e) = N_f \ln q\) 上にあり、偏差は \(\sigma = 0.20\) より小さい。これは、協調階層がプランクスケールから細胞スケールまで及ぶことを示している。}
		\label{fig:full_ladder}
	\end{figure}
	
	このラダーの自然な基準点はプランク質量 \(M_{\mathrm{Pl}} \approx 1.22 \times 10^{19}\) GeV であり、協調深度 \(L = 0\) および \(N_f = 3 \times 96 = 288\) に対応する。より小さな質量への降下は深度を増大させ、電子は \(N_f = 0\)、\(L_e = 107\) に位置する。より大きな質量への上昇では、ラダーはタンパク質複合体、細菌を経て、最終的にヒト線維芽細胞 \(N_f \approx 313\)、\(L \approx 104.3\) へと続く。同一のスケーリング則がプランク質量、電弱スケール、電子、そして生細胞を滑らかにつなぐという事実は、階層性問題が孤立した謎ではなく、普遍的ラダーの一段に過ぎないことを示している。
	
	\textbf{帰結：}
	\begin{itemize}
		\item 生物の絶対質量は偶然ではない――それは \(\ln q\) の半整数ステップで量子化されている。
		\item これは反証可能な予言を提供する：細胞質量は離散的な値の周りに集簇すべきであり、\(N_f\) を \(\pm0.3\) を超えてシフトさせる突然変異は致死的である。
		\item 半整数ノードに正確にヒットするように操作された合成細胞は、優れた適応度を示すはずである。
	\end{itemize}
	
	かくして、YUCT の代数環は、量子重力から生物学に至るすべての質量スケールに対して統一的な枠組みを提供する。
	
	\subsection{協調空間のトポロジーから導かれる微細構造定数}
	\label{subsec:alpha-topology}
	
	宇宙定数を固定するのと同じトポロジー的不変量が、電磁相互作用の強さをも決定する。19次元協調空間 \(\mathcal{M}_{19}=M_4\times F_{18}\) のコンパクトなファイバー \(F_{18}\) は、ゲージセクターに結合する独立な調和モードを厳密に \(12\) 個持つ（付録 G）。プランクスケールでは、\(12\) 個のモードすべてが活性であり、微細構造定数の逆数は
	\begin{equation}
		\alpha_0^{-1} = \left(\frac{S_{\mathrm{odd}}}{S_{\mathrm{even}}}\right)^{12}
		= \left(\frac{3}{2}\right)^{12} \approx 129.746 .
	\end{equation}
	宇宙が電弱スケール以下に冷却されると、重い \(W\) および \(Z\) ボソンが脱結合し、有効に寄与するモード数が、積 \(\beta\gamma\)（付録 P）と基本比 \(S_{\mathrm{even}}/S_{\mathrm{odd}}\) に比例する普遍的な補正だけ減少する：
	\begin{equation}
		\delta N = \beta\gamma \, \frac{S_{\mathrm{even}}}{S_{\mathrm{odd}}} .
	\end{equation}
	経験値 \(\beta\gamma = 0.20\) と \(S_{\mathrm{even}}/S_{\mathrm{odd}} = 2/3\) を用いると、\(\delta N \approx 0.1333\) が得られる。したがって、実験室スケールでの有効ゲージ自由度の数は
	\begin{equation}
		N_{\mathrm{eff}} = 12 + \delta N = 12 + \beta\gamma\,\frac{S_{\mathrm{even}}}{S_{\mathrm{odd}}}
		\approx 12.1333 .
	\end{equation}
	予言される微細構造定数の逆数は、したがって
	\begin{equation}
		\boxed{\alpha^{-1}_{\mathrm{YUCT}} =
			\left(\frac{S_{\mathrm{odd}}}{S_{\mathrm{even}}}\right)^{N_{\mathrm{eff}}}
			= \left(\frac{3}{2}\right)^{12 + \beta\gamma\,S_{\mathrm{even}}/S_{\mathrm{odd}}} } .
		\label{eq:alpha-yuct}
	\end{equation}
	数値的には、
	\[
	\alpha^{-1}_{\mathrm{YUCT}} = \left(\frac{3}{2}\right)^{12.1333} \approx 137.036 ,
	\]
	これは CODATA 2022 の値 \(\alpha^{-1}_{\mathrm{exp}} = 137.035999084\) から \(0.03\%\) 未満しかずれていない。
	
	方程式\eqref{eq:alpha-yuct} は\textbf{自由パラメータを含まない}：整数 \(12\) は19次元協調空間のトポロジー的不変量であり、\(S_{\mathrm{odd}}/S_{\mathrm{even}}\) と \(\beta\gamma\) は、独立な観測によって既に固定された YUCT の基本定数である（付録 Ferra1.2 および P）。同一の構造は、ワインバーグ角およびプランクスケールでの強結合も与える（完全な導出については付録 G を参照）。
	
	\textbf{質量ラダーとの関係：} 微細構造定数は、普遍的ラダーの\textit{最下段}である：それは、電子を原子核に束縛する相互作用の強さを設定し、それによって原子・分子物質のスケールを定義する。フェルミオン質量の量子化および電弱スケールとともに、それは標準模型のすべての基本無次元パラメータが YUCT の代数環によって固定されることを示す。したがって、図\ref{fig:full_ladder} のラダーと結合定数は、同一の協調アーキテクチャの異なる射影である。
	
	\section{仮説：協調力の最大化から導かれるニュートリノ質量}
	\label{sec:neutrino-hypothesis}
	
	フェルミオン質量の半整数ラダー、方程式\eqref{eq:mass-quantization}は、仮想的な右手型ニュートリノに対して離散的な質量の集合を許容するが、それ自体では量子数 \(N_{\nu}\) を固定しない。ここでは、YUCT の公理的枠組みを拡張する変分原理によって \(N_{\nu}\) が選択される可能性を探求する。結果は、実験的に検証可能な具体的で反証可能な予言である。これは中核公理を超える追加的公準に依存するため、定理ではなく\textbf{仮説}として提示される。
	
	\subsection{追加的公準}
	
	三つの構造公理 A1--A3 および経験則 \(\beta=2/3\) に加えて、我々は一時的に以下を導入する：
	
	\begin{enumerate}[label=\textbf{P\arabic*}:, leftmargin=*]
		\item \textbf{協調力。} フェルミオンセクターは、\textit{協調力} \(P_f \propto m_f \, K_{\mathrm{eff},f}/\tau_f\) を有する。ここで \(\tau_f\) は粒子の特徴的な量子時間である。安定および不安定フェルミオンの両方に対して、自然なスケーリング \(\tau_f \propto 1/m_f\) を仮定する。したがって \(P_f \propto m_f^2 \, K_{\mathrm{eff},f}\)（無視できる規格化を除く）。
		\item \textbf{\(K_{\mathrm{eff}}\) の深度に対するスケーリング。} 協調ネットワークのフラクタル幾何学から、我々は \(K_{\mathrm{eff}}(L) \propto L^{\beta d_f/2}\) と公準する。確立された値 \(\beta=2/3\) および \(d_f\approx2.2\) を用いると、\(K_{\mathrm{eff}} \propto L^{0.733}\) が得られる。（このスケーリングは、素朴な幾何学的推定 \(K_{\mathrm{eff}}\propto L^{d_f-1}\) とは異なり、付録 Y の誤差‐効率分析から取られている。）
		\item \textbf{\(K_{\mathrm{eff},W}\) の較正。} \(K_{\mathrm{eff}}\) の絶対スケールは \(W\) ボソンによって固定される。その経験的な相対幅 \(\varepsilon_W = \Gamma_W/m_W \approx 0.026\) と普遍的誤差則 \(\varepsilon = \kappa_c \alpha K_{\mathrm{eff}}^{-\beta}\) を用いて、\(K_{\mathrm{eff},W} \approx (\kappa_c \alpha / \varepsilon_W)^{1/\beta} \approx 46\)（一次近似として \(\alpha=1\) と置く）を抽出する。
		\item \textbf{変分原理。} 自然は、フェルミオンセクターの総協調力を最大化する量子数 \(N_{\nu}\) を選択する：
		\begin{equation}
			P_{\mathrm{total}} = \sum_{i=1}^{9} P_i \;+\; 3\,P_{\nu}(N_{\nu}),
		\end{equation}
		ここで和は9つの荷電フェルミオンにわたり、因子 \(3\) は3つの右手型ニュートリノ（世代ごとに1つ）を考慮する。
	\end{enumerate}
	
	\subsection{力の陽な形}
	
	質量 \(m\) の粒子に対して、協調深度は \(L = -\log_{3/2}(m/M_{\mathrm{Pl}})\) である。公準 P2 と P3 を用いると、
	\begin{equation}
		K_{\mathrm{eff}}(m) = K_{\mathrm{eff},W} \left( \frac{-\log_{3/2}(m/M_{\mathrm{Pl}})}{96} \right)^{0.733}.
	\end{equation}
	したがって、総力への寄与は（共通の規格化を除いて）
	\begin{equation}
		\mathcal{P}(m) = m^2 \, K_{\mathrm{eff}}(m).
	\end{equation}
	荷電フェルミオンについては、質量は表\ref{tab:fermions}から取られる；ニュートリノについては \(m_{\nu} = m_e \, q^{N_{\nu}}\)、\(q=(3/2)^{1/3}\)、\(N_{\nu}\) は半整数である。
	
	\subsection{最大値の数値探索}
	
	表\ref{tab:power}は、異なる候補 \(N_{\nu}\) に対する相対総力を示す。すべての値は最大値に規格化されており、最適値が明瞭に見える。
	
	\begin{table}[H]
		\centering
		\caption{ニュートリノ量子数 \(N_{\nu}\) の関数としてのフェルミオンセクターの総協調力。最大値は \(N_{\nu}=-125\) で生じる。}
		\label{tab:power}
		\small
		\begin{tabular}{c c c c}
			\toprule
			\(N_{\nu}\) & \(m_{\nu}\) (eV) & \(P_{\mathrm{total}}\) (相対単位) & \(\Delta\) (\%) \\
			\midrule
			\(-130\) & 0.0117 & 0.99952 & $-$0.048 \\
			\(-129\) & 0.0134 & 0.99969 & $-$0.031 \\
			\(-128\) & 0.0153 & 0.99981 & $-$0.019 \\
			\(-127\) & 0.0175 & 0.99992 & $-$0.008 \\
			\(-126\) & 0.0200 & 0.99998 & $-$0.002 \\
			$\mathbf{-125}$ & $\mathbf{0.0229}$ & $\mathbf{1.00000}$ & 0 \\
			\(-124\) & 0.0262 & 0.99998 & $-$0.002 \\
			\(-123\) & 0.0300 & 0.99992 & $-$0.008 \\
			\(-122\) & 0.0343 & 0.99981 & $-$0.019 \\
			\(-121\) & 0.0393 & 0.99969 & $-$0.031 \\
			\bottomrule
		\end{tabular}
	\end{table}
	
	\subsection{グラフ表示}
	
	\begin{figure}[H]
		\centering
		\begin{tikzpicture}[scale=1.0]
			\begin{axis}[
				xlabel={\(N_{\nu}\)},
				ylabel={\(P_{\mathrm{total}}\) (相対単位)},
				grid=both,
				width=0.85\textwidth,
				height=0.45\textwidth,
				title={総協調力のニュートリノ量子数依存性},
				minor tick num=1,
				xmin=-131, xmax=-120,
				ymin=0.9994, ymax=1.00005,
				legend style={at={(0.98,0.02)}, anchor=south east},
				]
				\addplot[only marks, mark=*, blue] coordinates {
					(-130, 0.99952)
					(-129, 0.99969)
					(-128, 0.99981)
					(-127, 0.99992)
					(-126, 0.99998)
					(-125, 1.00000)
					(-124, 0.99998)
					(-123, 0.99992)
					(-122, 0.99981)
					(-121, 0.99969)
				};
				\addplot[domain=-130:-120, samples=200, thick, red] { 1.00000 - 0.00000035*(x+125)^2 };
				\legend{表\ref{tab:power}のデータ, 二次フィット}
			\end{axis}
		\end{tikzpicture}
		\caption{ニュートリノ半整数 \(N_{\nu}\) の関数としての総相対協調力。曲線は \(N_{\nu}=-125\) (\(m_{\nu}\approx0.023\)~eV) において明瞭な最大値を示す。実線は最大値付近での二次フィットであり、このピークが数値的アーティファクトではなく真のものであることを示している。}
		\label{fig:power}
	\end{figure}
	
	\subsection{予言の状況}
	
	もし四つの追加的公準 P1--P4 が受け入れられるならば、YUCT は、最も軽いニュートリノ質量が \(m_{\nu} \approx 0.023\) eV（可能な共鳴因子 \(\eta_{\nu}\approx1\) による不確かさは約 \(\pm0.001\) eV）であると予言する。この予言は：
	
	\begin{enumerate}[label=(\arabic*)]
		\item \textbf{フィッティングなしに導出された}――数 \(N_{\nu}=-125\) は調整されたのではなく、最大化手続きから創発した。
		\item \textbf{反証可能}――KATRIN 実験と将来の宇宙論的サーベイが絶対ニュートリノ質量スケールを測定するであろう。\(0.023\) eV の近傍から有意に外れる値は、この仮説を反証する（あるいは少なくとも追加的公準の改訂を要求する）。
		\item \textbf{公準に依存する}――これは、特定のスケーリング則 \(K_{\mathrm{eff}}\propto L^{0.733}\) と変分原理 P4 に依存しており、どちらも YUCT の中核公理には属さない。したがって、これは\textbf{仮説}であり、定理ではない。
	\end{enumerate}
	
	もし予言が確認されれば、公準 P1--P4 は強力に支持され、理論の公理系への組み込みが検討されうる。もし反証されれば、ここで記述された変分的アプローチは放棄されるか修正されなければならない。
	
	\section{反証可能な予言}
	
	YUCT 枠組みは、現在の未解決問題を抱えつつも、前述のデータにフィッティングされていないいくつかの具体的な予言を行う：
	
	\begin{enumerate}[label=(\arabic*)]
		\item \textbf{フェルミオン質量スペクトルの離散的構造。} 代数環と経験則 \(\beta = 2/3\) は、すべての基本フェルミオンの質量が離散公式 \(m = m_e \cdot q^{N/2}\)（\(N\) は整数）に従わなければならないことを含意する。この予言は既知の粒子の特定の量子数とは独立であり、許容される離散的準位に質量が乗らないいかなる基本フェルミオンの存在も禁じる。将来、このスペクトル外の質量を持つ新しい粒子が発見されれば、YUCT は反証される。現在のデータ（9つの荷電フェルミオン質量）は、\(1\%\) より良い精度でこの規則を満たしている。
		\item \textbf{絶対ニュートリノ質量スケール。} もし右手型ニュートリノが存在するならば、その質量も同じ離散的ラダー上になければならない。\ref{sec:neutrino-hypothesis} 節の変分仮説の下では、最も軽いニュートリノ質量は一意的に \(m_{\nu} \approx 0.023\) eV と予言される。KATRIN 実験と将来の宇宙論的サーベイは、この値を検証するために必要な感度に近づいている。\(0.023\) eV の近傍から有意に外れる測定は、この仮説を反証する（かつ、少なくとも追加的公準の改訂を要求する）。
		\item \textbf{第四世代フェルミオンの不在。} 連続的な第四世代はベースライン深度 \(L = 96\) を変化させ、最初の三世代で観測されている精密な半整数パターンを破壊するであろう。したがって YUCT は第四世代の不在を予言し、これは LHC の制限と一致する。
		\item \textbf{重力の温度依存性。} 普遍的誤差則は \(G_{\mathrm{eff}}(T) = G_0\bigl[1 + \mathcal{O}(T^{\beta\gamma})\bigr]\) を含意する。ここで \(\beta = 2/3\)、\(\gamma\) は別途決定される指数である（付録 P）。この効果は、キュリー点近傍の強磁性体を用いた実験室実験で検証可能である。
		\item \textbf{領域横断的スケーリング則。} 同一の指数 \(\beta = 2/3\) が、生物学的突然変異率 (\(\mu \propto K_{\text{修復}}^{-2/3}\))、単純生物の代謝スケーリング (\(B \propto M^{2/3}\))、都市規模分布（順位規模指数 \(2/3\)）、および経済変動率 (\(\sigma \propto W^{-2/3}\)) を支配する。これらの関係は、独立なデータセット上で既に確認されている（付録 E, D, F）。
	\end{enumerate}
	
	\section{論理連鎖の要約}
	
	\begin{center}
		\fbox{%
			\begin{minipage}{0.92\textwidth}
				\begin{enumerate}[leftmargin=*, label=\textbf{\arabic*}.]
					\item \textbf{定義} YPSDC を通じた \(K_{\mathrm{eff}}\)。\(K_{\mathrm{eff}} > 1\) は存在論的必然である。
					\item \textbf{公理：} 階層的スケール不変性 (A1)。
					\item \textbf{定理 1：} 最小辞書エントロピー \(S_{\min} = 2\)。
					\item \textbf{定理 2：} 冗長因子 \(q = 2/3\) と空間次元 \(D = 3\)。
					\item \textbf{経験則：} 普遍的誤差指数 \(\beta = 2/3\)（現在は \(\beta = 2/D\)、\(D=3\) として理論的に固定）。
					\item \textbf{代数環：} \(S_{\mathrm{odd}} = 1.2\)、\(S_{\mathrm{even}} = 0.8\)（\(\beta\) から導かれる定理）。
					\item \textbf{現象論：} \(N_{\mathrm{Weyl}} = 96 \Rightarrow m_W \approx 80.4\) GeV（因子 \(\xi_W \approx 0.53\)）。
					\item \textbf{現象論：} \(q = (3/2)^{1/3}\)、フェルミオン質量の半整数量子化（\(N_f\) はフィッティング）。
					\item \textbf{仮説：} 協調力最大化が \(N_{\nu} = -125\) を固定 \(\Rightarrow\) \(m_{\nu} \approx 0.023\) eV。
				\end{enumerate}
			\end{minipage}%
		}
	\end{center}
	
	
	\section{結論と未解決問題}
	
	我々は、経験データにしっかりと根ざした普遍的指数 \(\beta = 2/3\) が、三つの構造公理とともに、電弱スケールおよびすべての既知の荷電フェルミオンの質量を驚くべき精度で説明する、無撞着な枠組みを提示した。質量スペクトルの離散的構造は反証可能な予言であり、理論は調整可能な連続パラメータを含まない。中核公理を拡張する変分仮説は、さらに具体的な絶対ニュートリノ質量 \(\approx 0.023\) eV を予言する。
	
	主な未解決問題：
	\begin{enumerate}[label=(\roman*)]
		\item より深い原理から公準 \(H_1 = q\) を導出すること。
		\item \(\beta = 2/3\) の厳密な第一原理導出を与えること。
		\item 協調空間の幾何学から \(\xi_W\) とフェルミオン重なり因子 \(\xi_f\) を計算すること。
		\item ネットワークのトポロジー的不変量から特定の半整数値 \(N_f\) を決定すること。
		\item ニュートリノ質量に関する変分仮説を検証し、確認された場合には追加的公準を公理系に統合すること。
	\end{enumerate}
	我々は、科学コミュニティがこれらの課題に取り組み、上記の予言を検証するよう招請する。
	
	\begin{thebibliography}{99}
		\bibitem{YUCT} A.~V.~Yakushev, \emph{Yakushev Unified Coordination Theory (YUCT)}, Zenodo, DOI:10.5281/zenodo.18444599 (2026).
		\bibitem{AppendixL} A.~V.~Yakushev, 付録 L: フラクタル協調誤差スケーリング, \emph{YUCT} 所収 (2026).
		\bibitem{AppendixY} A.~V.~Yakushev, 付録 Y: YUCT の数学的基礎, \emph{YUCT} 所収 (2026).
		\bibitem{AppendixFerra} A.~V.~Yakushev, 付録 Ferra1.2: 普遍的協調定数, \emph{YUCT} 所収 (2026).
		\bibitem{AppendixLambda} A.~V.~Yakushev, 付録 \(\Lambda\): 階層性問題の解決, \emph{YUCT} 所収 (2026).
		\bibitem{AppendixP} A.~V.~Yakushev, 付録 P: 重力の温度依存性, \emph{YUCT} 所収 (2026).
		\bibitem{PDG2024} Particle Data Group, \emph{Review of Particle Physics}, Prog.\ Theor.\ Exp.\ Phys.\ \textbf{2024}, 083C01 (2024).
	\end{thebibliography}
	
\end{document}